% مجموعهٔ تمرین‌ها — درس ۲
% شمارهٔ درس از نام این پرونده تعیین می‌شود.

\begin{exo}[از یخ سرد تا آب جوشان: مرتبه‌های بزرگی]{chauffage-glace-eau-ebullition}
\seoready{true}
\keywords{گرماسنجی, تغییر فاز, ذوب, تبخیر, گرمای نهان, مرتبهٔ بزرگی}

\textbf{داده‌های عددی (در فشار جو):}
\[
\begin{aligned}
c_{\mathrm{i}} &= 2{,}1\times10^3\,\text{J}\!\cdot\!\text{kg}^{-1}\!\cdot\!\text{K}^{-1},
& L_f &= 3{,}34\times10^5\,\text{J}\!\cdot\!\text{kg}^{-1},\\
c_{\mathrm{w}} &= 4{,}18\times10^3\,\text{J}\!\cdot\!\text{kg}^{-1}\!\cdot\!\text{K}^{-1},
& L_v &= 2{,}26\times10^6\,\text{J}\!\cdot\!\text{kg}^{-1}.
\end{aligned}
\]
در فشار جو، یک کیلوگرم یخ را که در آغاز دمای $-100\,{}^\circ\text{C}$ دارد به‌تدریج گرم می‌کنیم.

\begin{enumerate}
  \item انرژی لازم برای رساندن یخ از $-100\,{}^\circ\text{C}$ به $0\,{}^\circ\text{C}$ را محاسبه کنید.
  \item انرژی لازم برای ذوب کامل یخ در $0\,{}^\circ\text{C}$ را محاسبه کنید.
  \item انرژی لازم برای گرم‌کردن آب مایع از $0\,{}^\circ\text{C}$ تا $100\,{}^\circ\text{C}$ را محاسبه کنید.
  \item انرژی اضافی لازم برای تبخیر کامل آب در $100\,{}^\circ\text{C}$ را محاسبه کنید.
  \item کدام مرحله بیشترین انرژی را می‌خواهد؟ انرژی گرم‌کردن یک کیلوگرم آب از $0\,{}^\circ\text{C}$ تا $100\,{}^\circ\text{C}$ را با انرژی پتانسیل جرم $m$ که از ارتفاع ده متر سقوط می‌کند مقایسه کنید. برای برابری این انرژی‌ها چه جرمی باید سقوط کند؟
  \item اکنون فرایند وارون را در نظر بگیرید. جرم $m_v$ از بخار آب روی شیشهٔ جلوی خودرو چگالیده می‌شود، بی‌آنکه آب مایع حاصل سپس سردتر شود. گرمای آزادشدهٔ $Q_{\mathrm{rel}}$ در چگالش را بنویسید و آن را برای $m_v=1{,}0\times10^{-2}\,\text{kg}$ محاسبه کنید. در دمای شیشه $L_v=2{,}45\times10^6\,\text{J}\!\cdot\!\text{kg}^{-1}$ بگیرید.
\end{enumerate}

\begin{indication}
برای گرم‌کردن یک فاز بدون تغییر فاز از $Q=mc\Delta T$ و برای تغییر فاز در دمای ثابت از $Q=mL$ استفاده کنید. انرژی پتانسیل گرانشی جرم $m$ در ارتفاع $h$ برابر $E_p=mgh$ است؛ $g=9{,}81\,\text{m}\!\cdot\!\text{s}^{-2}$ بگیرید.
\end{indication}

\begin{solution}
\textbf{1.} گرم‌کردن یخ تا دمای ذوب نیازمند
\[
Q_1=mc_{\mathrm{i}}\Delta T
=1{,}0\times2{,}1\times10^3\times100
=2{,}1\times10^5\,\text{J}
\]
است.

\textbf{2.} هنگام ذوب، دما روی $0\,{}^\circ\text{C}$ می‌ماند:
\[
Q_2=mL_f=1{,}0\times3{,}34\times10^5=3{,}34\times10^5\,\text{J}.
\]

\textbf{3.} برای گرم‌کردن آب مایع از $0$ تا $100\,{}^\circ\text{C}$،
\[
Q_3=mc_{\mathrm{w}}\Delta T
=1{,}0\times4{,}18\times10^3\times100
=4{,}18\times10^5\,\text{J}.
\]

\textbf{4.} تبخیر کامل به
\[
Q_4=mL_v=1{,}0\times2{,}26\times10^6=2{,}26\times10^6\,\text{J}
\]
دیگر نیاز دارد. همین مرحله به‌تنهایی انرژی‌ای در حد چند مگاژول می‌خواهد.

\textbf{5.} تبخیر با فاصلهٔ زیاد پرهزینه‌ترین مرحله است:
\[
Q_4=2{,}26\times10^6\,\text{J}
>Q_3=4{,}18\times10^5\,\text{J}
>Q_2=3{,}34\times10^5\,\text{J}
>Q_1=2{,}1\times10^5\,\text{J}.
\]
تبخیر تقریباً $(2{,}26\times10^6)/(4{,}18\times10^5)\approx5{,}4$ برابر گرم‌کردن آب مایع از $0$ تا $100\,{}^\circ\text{C}$ انرژی می‌خواهد.

برای جرم $m$ که از $h=10\,\text{m}$ سقوط می‌کند، $E_p=mgh$ است. با قرار دادن $mgh=Q_3$،
\[
m=\frac{Q_3}{gh}
=\frac{4{,}18\times10^5}{9{,}81\times10}
\approx4{,}26\times10^3\,\text{kg}.
\]
پس باید جرمی در حدود $4{,}3$ تن را از ارتفاع ده متر رها کرد تا به‌اندازهٔ گرم‌کردن یک کیلوگرم آب از $0$ تا $100\,{}^\circ\text{C}$ انرژی آزاد شود! ازاین‌رو گرم‌کردن آب سهم مهمی از مصرف انرژی خانگی دارد. ظرفیت گرمایی ویژهٔ آب نیز در قیاس با فلزهای معمول بسیار زیاد است؛ برای نمونه تقریباً ده برابر مس (تمرین‌های بعدی را ببینید).

\textbf{6.} چگالش وارون تبخیر است: بخار گرمای نهان را به شیشه می‌دهد،
\[
Q_{\mathrm{rel}}=m_vL_v.
\]
برای $m_v=1{,}0\times10^{-2}\,\text{kg}$،
\[
Q_{\mathrm{rel}}=1{,}0\times10^{-2}\times2{,}45\times10^6
=2{,}45\times10^4\,\text{J}.
\]
بنابراین چگالش جرم اندکی بخار نیز می‌تواند گرمای چشمگیری آزاد کند. شیشه و محیط پیرامونش این گرما را می‌گیرند.
\end{solution}
\end{exo}

\begin{exo}[تبخیر و تعرق یک درخت بزرگ: تهویهٔ مطبوع طبیعی؟]{odg-evapotranspiration-arbre}
\seoready{true}
\keywords{مرتبهٔ بزرگی, تبخیر و تعرق, درخت, گرمای نهان, توان سرمایشی, تهویهٔ مطبوع, COP}

در روزی گرم و خشک، یک درخت بزرگ که آب کافی در اختیار دارد می‌تواند حدود $500\,\text{L}=5{,}0\times10^{-1}\,\text{m}^3$ آب را با تبخیر و تعرق از دست بدهد. چون تعرق عمدتاً در روشنایی رخ می‌دهد، فرض می‌کنیم طی دوازده ساعت پخش شده است. تاج درخت روی زمین مساحت $A_{\mathrm{tree}}=50\,\text{m}^2$ را می‌پوشاند. چگالی آب و گرمای نهان ویژهٔ تبخیر آن در فشار جو برابرند با
\[
\rho_{\mathrm{w}}=1{,}0\times10^3\,\text{kg}\!\cdot\!\text{m}^{-3},
\qquad L_v=2{,}45\times10^6\,\text{J}\!\cdot\!\text{kg}^{-1}.
\]

\begin{enumerate}
  \item جرم آبی را که درخت طی روز تبخیر و تعرق می‌کند و سپس انرژی لازم برای تبخیر آن را محاسبه کنید.
  \item توضیح دهید چرا این فرایند اثر سرمایشی دارد.
  \item توان سرمایشی میانگین درخت به‌ازای هر متر مربع را طی دوازده ساعت تعرق فعال محاسبه کنید.
  \item برای مقایسه، دستگاه تهویهٔ مطبوعی با توان الکتریکی $P_{\mathrm{elec}}=2{,}0\times10^3\,\text{W}$ و ضریب عملکرد $\mathrm{COP}=3{,}0$ اتاقی با مساحت $A_{\mathrm{room}}=20\,\text{m}^2$ را خنک می‌کند. بنا بر تعریف،
\[
\mathrm{COP}=\frac{P_{\mathrm{cool}}}{P_{\mathrm{elec}}},
\]
که $P_{\mathrm{cool}}$ توان گرمایی خارج‌شده از اتاق است. $P_{\mathrm{cool}}$ و توان سرمایشی به‌ازای هر متر مربع اتاق را محاسبه و با درخت مقایسه کنید.
  \item چرا حتی اگر توان بر واحد سطح آن‌ها هم‌مرتبه باشد نمی‌توان نتیجه گرفت درخت و دستگاه تهویه دقیقاً سرمایش احساس‌شدهٔ یکسانی ایجاد می‌کنند؟
  \item به نظر شما آیا می‌توان با تعداد زیادی گیاه آپارتمانی خانه را خنک کرد؟ (دانش عمومی: پاسخ به فرایندهای زیستی وابسته است.)
\end{enumerate}

\begin{indication}
از $m=\rho V$، $Q_{\mathrm{evap}}=mL_v$ و $P=Q/\tau$ استفاده کنید. برای دستگاه تهویه، $P_{\mathrm{cool}}=\mathrm{COP}\,P_{\mathrm{elec}}$.
\end{indication}

\begin{solution}
\textbf{1.} $500\,\text{L}=5{,}0\times10^{-1}\,\text{m}^3$، پس
\[
m=\rho_{\mathrm{w}}V=1{,}0\times10^3\times5{,}0\times10^{-1}=5{,}0\times10^2\,\text{kg},
\]
و
\[
Q_{\mathrm{evap}}=mL_v=5{,}0\times10^2\times2{,}45\times10^6=1{,}225\times10^9\,\text{J}.
\]
مرتبهٔ بزرگی آن یک گیگاژول است.

\textbf{2.} تبخیر تغییری گرماگیر است. انرژی $Q_{\mathrm{evap}}$ از برگ‌ها، خاک و هوای اطراف گرفته می‌شود و با کاهش انرژی گرمایی، دمای آن‌ها پایین می‌آید.

\textbf{3.} دوازده ساعت برابر $\tau=12\times3600=4{,}32\times10^4\,\text{s}$ است. بنابراین
\[
P_{\mathrm{tree}}=\frac{Q_{\mathrm{evap}}}{\tau}
=\frac{1{,}225\times10^9}{4{,}32\times10^4}
\approx2{,}84\times10^4\,\text{W},
\]
و روی $50\,\text{m}^2$،
\[
\frac{P_{\mathrm{tree}}}{A_{\mathrm{tree}}}\approx5{,}7\times10^2\,\text{W}\!\cdot\!\text{m}^{-2}.
\]

\textbf{4.}
\[
P_{\mathrm{cool}}=\mathrm{COP}\,P_{\mathrm{elec}}=3{,}0\times2{,}0\times10^3=6{,}0\times10^3\,\text{W},
\]
و
\[
\frac{P_{\mathrm{cool}}}{A_{\mathrm{room}}}=\frac{6{,}0\times10^3}{20}=3{,}0\times10^2\,\text{W}\!\cdot\!\text{m}^{-2}.
\]
در این مدل، تبخیر و تعرق درخت تقریباً دو برابر یک دستگاه تهویهٔ معمولی توان دارد.

\textbf{5.} این مقایسه تنها جریان‌های انرژی را می‌سنجد. دستگاه تهویه حجمی بسته و کنترل‌شده را خنک می‌کند، اما بخار درخت در جو پخش می‌شود و باد هوای گرم بازمی‌آورد. اثر احساس‌شده به تهویه، رطوبت، دمای محیط و فاصله از درخت وابسته است. تعرق نیز با نور خورشید، باد، رطوبت، گونه و آب در دسترس تغییر می‌کند و در خشکسالی ممکن است بسیار کم شود. از سوی دیگر، درخت سایه هم فراهم می‌کند. پس محاسبه‌ها مقایسهٔ مرتبهٔ بزرگی‌اند، نه برابری دقیق آسایش گرمایی.

\textbf{6.} در عمل اثر سرمایشی گیاهان آپارتمانی بسیار کم است. در بیشتر گیاهان نور روزنه‌های برگ را باز می‌کند تا بخار آب خارج شود؛ نور داخل ساختمان معمولاً کم است و روزنه‌ها کمتر باز می‌شوند. در اتاقی با تهویهٔ کم، افزایش رطوبت نیز تبخیر را به‌تدریج کند می‌کند. تعداد زیادی گیاه شاید سرمایش تبخیری موضعی اندکی ایجاد کنند، اما جای دستگاه تهویه را نمی‌گیرند. برای دفع پیوستهٔ انرژی باید هوای مرطوب به بیرون رانده شود. نور مصنوعی شدید شاید تعرق را بالا ببرد، ولی انرژی الکتریکی آن سرانجام تقریباً به‌طور کامل در اتاق به گرما تبدیل می‌شود. گیاهان CAM سازگار با محیط خشک، مانند کاکتوس‌ها، استثنا هستند: روزنه‌هایشان عمدتاً شب باز می‌شود.
\end{solution}
\end{exo}

\begin{exo}[مخلوط‌کردن دو جرم آب]{calorimetrie-melange-eau}
\seoready{true}
\keywords{گرماسنجی, مخلوط, ظرفیت گرمایی, ظرفیت گرمایی ویژه, جوزف بلک, دمای تعادل}

جرم $m_1=0{,}200\,\text{kg}$ آب با دمای $T_1=80\,{}^\circ\text{C}$ را در ظرفی می‌ریزیم که از پیش $m_2=0{,}300\,\text{kg}$ آب با دمای $T_2=15\,{}^\circ\text{C}$ دارد. ظرف یک گرماسنج ایده‌آل است: اتلاف گرما به بیرون و ظرفیت گرمایی خود ظرف را نادیده می‌گیریم. رابطهٔ
\[
Q=mc\Delta T
\]
را به یاد آورید که در آن ظرفیت گرمایی ویژهٔ آب مایع $c=4{,}18\times10^3\,\text{J}\!\cdot\!\text{kg}^{-1}\!\cdot\!\text{K}^{-1}$ است.

\begin{enumerate}
  \item نشان دهید گرمای ازدست‌رفتهٔ آب گرم به‌طور کامل به آب سرد می‌رسد؛ معادلهٔ پایستگی را برحسب $m_1$، $m_2$، $c$، $T_1$، $T_2$ و دمای تعادل $T_{eq}$ بنویسید.
  \item عبارت جبری $T_{eq}$ و سپس مقدار عددی آن را بیابید.
  \item گرمای $Q$ داده‌شده از سوی آب گرم هنگام مخلوط‌شدن را محاسبه کنید.
  \item آیا چنین مخلوطی از یک ماده می‌تواند مقدار $c$ را تعیین کند؟ توضیح دهید.
\end{enumerate}

\begin{indication}
چون گرماسنج عایق و صلب است، انرژی درونی کل $U_1+U_2$ پایسته می‌ماند: گرمایی که یکی می‌دهد با علامت مخالف به دیگری می‌رسد.
\end{indication}

\begin{solution}
\textbf{1.} در گرماسنج عایق، انرژی درونی کل سامانهٔ آب گرم و آب سرد پایسته است. با کاربرد $Q=mc\Delta T$ برای هر جرم از حالت آغازین تا تعادل،
\[
m_1c(T_{eq}-T_1)+m_2c(T_{eq}-T_2)=0.
\]

\textbf{2.} چون هر دو یک ماده‌اند، $c$ ساده می‌شود:
\[
T_{eq}=\frac{m_1T_1+m_2T_2}{m_1+m_2}
=\frac{0{,}200\times80+0{,}300\times15}{0{,}200+0{,}300}
=41\,{}^\circ\text{C}.
\]

\textbf{3.}
\[
Q=m_1c(T_1-T_{eq})
=0{,}200\times4{,}18\times10^3\times(80-41)
\approx3{,}26\times10^4\,\text{J}.
\]
آب سرد دقیقاً همین مقدار را می‌گیرد:
\[
m_2c(T_{eq}-T_2)=0{,}300\times4{,}18\times10^3\times26\approx3{,}26\times10^4\,\text{J}.
\]

\textbf{4.} خیر. برای دو جرم از یک ماده، همان $c$ هر دو جملهٔ موازنه را ضرب می‌کند و ساده می‌شود. دمای تعادل اندازه‌گیری‌شده به $c$ وابسته نیست و تنها میانگین دماهای آغازین با وزن جرم است. برای اندازه‌گیری ظرفیت گرمایی ویژه با روش اختلاط باید ماده‌ای با ظرفیت گرمایی معلوم، مانند تمرین بعدی، وارد کرد.
\end{solution}
\end{exo}

\begin{exo}[تعیین ظرفیت گرمایی ویژهٔ فلز با روش اختلاط]{calorimetrie-methode-melanges}
\seoready{true}
\keywords{گرماسنجی, روش اختلاط, ظرفیت گرمایی ویژه, گرماسنج, گرمای ویژه, معادل آبی}

می‌خواهیم مانند جوزف بلک در سدهٔ XVIII\textsuperscript{e}، ظرفیت گرمایی ویژهٔ $c_m$ یک فلز ناشناخته را با \emph{روش اختلاط} اندازه بگیریم. نمونه‌ای با جرم $m=0{,}150\,\text{kg}$ را تا $T_m=95\,{}^\circ\text{C}$ گرم و سریع در گرماسنجی فرو می‌بریم که $m_e=0{,}250\,\text{kg}$ آب در $T_e=18\,{}^\circ\text{C}$ دارد؛ $c_e=4{,}18\times10^3\,\text{J}\!\cdot\!\text{kg}^{-1}\!\cdot\!\text{K}^{-1}$. پس از رسیدن به تعادل گرمایی، $T_{eq}=22\,{}^\circ\text{C}$ اندازه‌گیری می‌شود. ابتدا ظرفیت گرمایی خود ظرف گرماسنج (دیواره‌ها، همزن و دماسنج) را نادیده بگیرید.

\begin{enumerate}
  \item مانند تمرین پیشین، پایستگی انرژی سامانهٔ عایق فلز و آب را طوری بنویسید که $c_m$ تنها مجهول باشد.
  \item عبارت جبری و مقدار عددی $c_m$ را بیابید. این مقدار تقریباً با کدام فلز معمول سازگار است؟ مقدارهای مرجع: آلومینیوم $\approx9{,}0\times10^2\,\text{J}\!\cdot\!\text{kg}^{-1}\!\cdot\!\text{K}^{-1}$، آهن $\approx4{,}5\times10^2$، مس $\approx3{,}9\times10^2$، سرب $\approx1{,}3\times10^2$.
  \item در واقع خود ظرف نیز بخشی از گرمای فلز را می‌گیرد. این اثر را با \emph{معادل آبی} $\mu=1{,}5\times10^{-2}\,\text{kg}$ مدل می‌کنیم: گرماسنج چنان رفتار می‌کند که گویی جرم اضافی $\mu$ از آب است. $c_m$ را با درنظرگرفتن $\mu$ دوباره حساب کنید و جهت تصحیح را توضیح دهید (آیا $c_m$ تصحیح‌شده از پاسخ ۲ بزرگ‌تر است یا کوچک‌تر؟).
  \item چرا از نظر تجربی لازم است نمونه را \emph{سریع} در گرماسنج فرو ببریم و هنگام اندازه‌گیری آن را از هوای محیط به‌خوبی عایق گرمایی کنیم؟
\end{enumerate}

\begin{indication}
فلز گرمای $Q_m=mc_m(T_m-T_{eq})$ را می‌دهد که به‌طور کامل به آب (و در پرسش ۳ به گرماسنجِ معادل جرم اضافی $\mu$ آب) می‌رسد: $Q_m=(m_e+\mu)c_e(T_{eq}-T_e)$.
\end{indication}

\begin{solution}
\textbf{1.} در سامانهٔ عایق، گرمایی که فلز با سردشدن از $T_m$ تا $T_{eq}$ می‌دهد به‌طور کامل به آبی می‌رسد که از $T_e$ تا $T_{eq}$ گرم می‌شود:
\[
mc_m(T_m-T_{eq})=m_ec_e(T_{eq}-T_e).
\]

\textbf{2.}
\[
c_m=\frac{m_ec_e(T_{eq}-T_e)}{m(T_m-T_{eq})}
=\frac{0{,}250\times4{,}18\times10^3\times(22-18)}{0{,}150\times(95-22)}
\approx3{,}82\times10^2\,\text{J}\!\cdot\!\text{kg}^{-1}\!\cdot\!\text{K}^{-1}.
\]
این مقدار به مقدار مس ($\approx3{,}9\times10^2\,\text{J}\!\cdot\!\text{kg}^{-1}\!\cdot\!\text{K}^{-1}$) بسیار نزدیک است.

\textbf{3.} اکنون آب و ظرف با هم گرما را می‌گیرند:
\[
c_m=\frac{(m_e+\mu)c_e(T_{eq}-T_e)}{m(T_m-T_{eq})}
=\frac{(0{,}250+0{,}015)\times4{,}18\times10^3\times4}{0{,}150\times73}
\approx4{,}05\times10^2\,\text{J}\!\cdot\!\text{kg}^{-1}\!\cdot\!\text{K}^{-1}.
\]
تصحیح $c_m$ را اندکی افزایش می‌دهد. با نادیده‌گرفتن معادل آبی در پرسش ۲، گرمای واقعی داده‌شده از سوی فلز—که بخشی از آن ظرف را گرم کرده بود—و در نتیجه $c_m$ را کمتر برآورد کرده بودیم.

\textbf{4.} روش بر فرض عایق‌بودن سامانه هنگام اندازه‌گیری استوار است. فروبردن سریع نمونه زمان تبادل گرما با هوا پیش از رسیدن به آب را کم می‌کند؛ عایق‌کردن خوب نیز اتلاف به بیرون را هنگام برقراری تعادل کاهش می‌دهد. هر نشت گرمایی حساب‌نشده موازنهٔ پرسش ۱ و مقدار $c_m$ حاصل را نادرست می‌کند. بلک و سپس لاووازیه و لاپلاس نیز در گرماسنج یخی خود همین دقت تجربی را به کار می‌بردند.
\end{solution}
\end{exo}

\begin{exo}[کالری‌های غذایی و توان سوخت‌وسازی]{calorimetrie-calories-alimentaires}
\seoready{true}
\keywords{کالری, کیلوکالری, کالری غذایی, معادل مکانیکی, توان سوخت‌وسازی, یکاها}

برچسب تغذیه نشان می‌دهد یک بزرگسال به‌طور میانگین حدود $2000\,\text{kcal}$ در روز مصرف می‌کند. کیلوکالری یکای غیر SI است که هنوز در تغذیه به کار می‌رود؛ تبدیل آن $1\,\text{kcal}=4{,}18\times10^3\,\text{J}$ است.

\begin{enumerate}
  \item این سهم انرژی روزانه را به ژول تبدیل کنید.
  \item با فرض مصرف یکنواخت انرژی طی ۲۴ ساعت، توان سوخت‌وسازی میانگین این بزرگسال را به وات به دست آورید.
  \item روی یک میلهٔ انرژی‌زا نوشته شده است «$210\,\text{kcal}$». اگر همین انرژی به‌طور کامل به گرما تبدیل شود، چه جرم آبی را می‌توان از $20\,{}^\circ\text{C}$ به نقطهٔ جوش ($100\,{}^\circ\text{C}$) رساند؟ $c_{\mathrm{w}}=4{,}18\times10^3\,\text{J}\!\cdot\!\text{kg}^{-1}\!\cdot\!\text{K}^{-1}$ بگیرید.
  \item به نظر شما این انرژی در بدن انسان واقعاً به چه صورت‌هایی مصرف می‌شود؟ (پاسخ کیفی)
\end{enumerate}

\begin{indication}
برای پرسش ۲ از $\mathcal P=E/\Delta t$ استفاده کنید. در پرسش ۳ ابتدا انرژی برچسب را به ژول تبدیل و سپس $m$ را از $Q=mc\Delta T$ جدا کنید.
\end{indication}

\begin{solution}
\textbf{1.} در یکاهای SI،
\[
E=2000\times4{,}18\times10^3=8{,}36\times10^6\,\text{J}.
\]

\textbf{2.}
\[
\mathcal P=\frac{8{,}36\times10^6}{8{,}64\times10^4}\approx97\,\text{W}.
\]
انسان در سراسر روز به‌طور میانگین توانی تقریباً برابر یک لامپ رشته‌ای مصرف می‌کند. این مرتبهٔ بزرگی شگفت‌آور سودمندی تبدیل کالری غذایی به یکاهای متداول فیزیکی را نشان می‌دهد.

\textbf{3.} انرژی میله برابر است با
\[
Q=210\times4{,}18\times10^3\approx8{,}78\times10^5\,\text{J}.
\]
با $\Delta T=80\,\text{K}$،
\[
m=\frac{Q}{c\Delta T}
=\frac{8{,}78\times10^5}{4{,}18\times10^3\times80}
\approx2{,}63\,\text{kg}.
\]
پس از نظر مرتبهٔ بزرگی، یک میلهٔ انرژی‌زا برای رساندن حدود $2{,}5\,\text{kg}$ آب شیر به نقطهٔ جوش انرژی دارد.

\textbf{4.} بدن انرژی را به معنای ناپدیدکردن آن «مصرف» نمی‌کند؛ انرژی شیمیایی مواد مغذی را تبدیل می‌کند. بخشی برای کار مکانیکی ماهیچه‌ها، کارکرد اندام‌ها، نگهداری شیب‌های الکتریکی سلول‌ها، جابه‌جایی مولکول‌ها و سنتزهای شیمیایی لازم برای نوسازی بافت‌ها به کار می‌رود. بخشی دیگر ممکن است به صورت شیمیایی، به‌ویژه در ذخایر گلیکوژن و چربی، انباشته شود. سرانجام بیشتر انرژی به‌کاررفته به شکل گرما پخش می‌شود و پیش از انتقال به محیط به حفظ دمای بدن کمک می‌کند. بخش کوچکی نیز به صورت انرژی شیمیایی باقی‌مانده در مواد دفعی از بدن بیرون می‌رود.
\end{solution}
\end{exo}
